% ID: SEL-D0009-BC-E06
% ESTADO: validada. S010; tres pasadas del mismo autor y comprobación matemática; sin segundo revisor.
\textbf{Sucesos y probabilidades del árbol.}\par\nopagebreak[4]
Sean $A,B$ los dos tipos de bacteria y $R$ la reacción al antibiótico. Son tipos exhaustivos y disjuntos: $P(A)=0.70$, $P(B)=0.30$. Los datos son condicionados: $P(R\mid A)=0.15$ y $P(R\mid B)=0.80$.

\textbf{a) Probabilidad de reacción.}\par\nopagebreak[4]
Aplicamos la probabilidad total:
\[
\boxed{P(R)=0.70\cdot0.15+0.30\cdot0.80=0.105+0.240=0.345.}
\]
\textbf{b) Tipo B, sabiendo que hay reacción.}\par\nopagebreak[4]
Bayes invierte la condición usando la probabilidad total anterior:
\[
\boxed{P(B\mid R)=\frac{P(B)P(R\mid B)}{P(R)}
=\frac{0.240}{0.345}=\frac{16}{23}\approx0.695652.}
\]
No debe confundirse este resultado con $P(R\mid B)=0.80$.

\textbf{c) Tipo A y sin reacción.}\par\nopagebreak[4]
Se pide una intersección, no una probabilidad condicionada:
\[
\boxed{P(A\cap R^c)=P(A)\bigl[1-P(R\mid A)\bigr]
=0.70\cdot0.85=0.595.}
\]
Comprobación: las cuatro hojas del árbol son $0.105$, $0.595$, $0.240$ y $0.060$; suman $1$ y las de reacción suman $0.345$.
